\documentclass[10pt]{article}
\usepackage[letterpaper,margin=0.75in]{geometry}
\usepackage{amsmath,amssymb,amsthm,booktabs}
\usepackage[T1]{fontenc}
\usepackage{lmodern}
\usepackage[colorlinks=true,urlcolor=blue,linkcolor=blue]{hyperref}
\usepackage{microtype}
\setlength{\parindent}{0pt}
\setlength{\parskip}{4pt}
\setlength{\emergencystretch}{2em}
\newtheorem{theorem}{Theorem}
\newcommand{\N}{\mathbb N}
\newcommand{\R}{\mathbb R}
\newcommand{\code}[1]{\texttt{\small #1}}
\title{The containment minimum for paths\\[3pt]
\large A disproof of Conjecture 00000002187 as written}
\author{C0ldSmi1e}
\date{9 October 2026}
\begin{document}
\maketitle
\vspace{-1.1em}

\begin{abstract}
The conjecture explicitly defines its invariant as the minimum number of
edges in a graph containing a path. Under that definition, the minimum for
an $n$-vertex path is exactly $n-1$ for $n\geq1$. Its ratio to $n$ tends to
$1$, contradicting the asserted coefficient $3$. The submitted Lean 4
project proves the attained minimum, its exact value, the actual limit,
and the full negation of the proposed little-$o$ formula.
\end{abstract}

\section{Statement and scope}
The English and Chinese originals in \code{ORIGINAL.md} both use ordinary
containment in their explicit definition. We therefore write
\[
 \rho(H)=\min\bigl\{|E(G)|: G\text{ contains a copy of }H,\ |E(G)|<\infty\bigr\}.
\]
Graphs are undirected and simple. A copy is an injective map of vertices
preserving adjacency; nonadjacency need not be preserved. The path $P_n$
has vertex set $\{0,\ldots,n-1\}$ and consecutive vertices are adjacent.
The claimed numerical assertion is $\rho(P_n)=(3+o(1))n$ as $n\to\infty$.

The source calls this invariant a size-Ramsey number. The conventional
definition requires a monochromatic copy in every two-colouring of the
host's edges~\cite{cfw}. That requirement is absent from the printed
definition. This submission follows the printed definition and does not
assert a result for that different, colouring-based invariant.
The sentence about a Bhamidi--Hoagy--Hoarau network limit specifies no
network family or limit and is presented as an explanation, not a
hypothesis restricting the numerical assertion. No such missing
objects are introduced here.

\section{Exact minimum and disproof}
\begin{theorem}
For every simple graph $H$ with finitely many edges, the above minimum
exists, is attained, and satisfies $\rho(H)=|E(H)|$.
\end{theorem}
\begin{proof}
The host $G=H$, with the identity copy, is admissible. If $f$ embeds a copy
of $H$ into an admissible $G$, it sends each unordered edge $\{u,v\}$ to
$\{f(u),f(v)\}$. Injectivity of $f$ implies injectivity of this edge map.
Consequently $|E(H)|\leq |E(G)|$ for every admissible host, while $H$
itself attains equality. This proves both minimum semantics and its value.
\end{proof}

For $n=k+1$, the edges of $P_n$ are precisely
$\{i,i+1\}$ for $0\leq i<k$. They are distinct, so $|E(P_n)|=k$.
Thus $\rho(P_n)=n-1$ for every $n\geq1$, and
\[
 \frac{\rho(P_n)}{n}=1-\frac1n\longrightarrow1.
\]
If a real sequence $\varepsilon_n\to0$ satisfied
$\rho(P_n)=(3+\varepsilon_n)n$ for all sufficiently large $n$, division
by positive $n$ would make the same ratio tend to $3$.
Uniqueness of real limits gives the contradiction $1=3$.

If the index instead counts edges, the path is $P_{n+1}$ in our notation
and its minimum is exactly $n$. The normalized coefficient is again $1$.
The discrepancy therefore does not depend on this indexing convention.

\newpage
\section{Correspondence with the formal theorem}
The mathematical source is \code{Conjecture2187.lean}. It uses Mathlib's
\code{SimpleGraph}, unordered edge sets, and \code{pathGraph n} on
\code{Fin n}. All the following names have prefix \code{Conjecture2187.}

\begingroup
\small
\renewcommand{\arraystretch}{1.10}
\begin{tabular}{@{}p{0.45\textwidth}p{0.51\textwidth}@{}}
\toprule
\normalfont Formal name & Meaning or established result\\
\midrule
\code{Contains} & An injective graph homomorphism; ordinary subgraph containment.\\
\code{Attainable} & An existential host graph with a finite edge set, a copy, and the specified edge count.\\
\code{minimumEdges} & The least attained natural edge count, defined using \code{Nat.find} and a proved existence statement.\\
\code{minimumEdges\_attained} & A genuine host attains the minimum.\\
\code{minimumEdges\_le} & The minimum is no larger than any admissible edge count.\\
\code{minimumEdges\_eq} & The minimum equals the contained graph's edge count.\\
\code{pathMinimum\_eq} & $\rho(P_n)=n-1$.\\
\code{pathMinimum\_succ} & $\rho(P_{n+1})=n$.\\
\code{pathMinimum\_ratio\_tendsto\_one} & $\rho(P_n)/n\to1$.\\
\code{conjecture\_false} & The full logical negation of \code{ClaimedAsymptotic}.\\
\bottomrule
\end{tabular}
\endgroup

The proposition \code{ClaimedAsymptotic} is exactly
\[
 \exists\varepsilon:\N\to\R,\quad
 \varepsilon_n\longrightarrow0
 \quad\land\quad
 \forall^{\text{eventually}} n\in\N,\quad
 (\rho(P_n):\R)=(3+\varepsilon_n)(n:\R).
\]
The final theorem proves its negation. The proof works with the weaker,
eventual equality, so it also excludes a reading requiring the equality
for every index. Ratio calculations are restricted eventually to
$n\geq1$; no conclusion relies on division by zero.

The extremum ranges over arbitrary host vertex sets in Lean's \code{Type},
with finite edge sets. Hosts with infinitely many isolated vertices are
allowed. The lower bound holds for every such host, and the finite path
attains it. Hence the result also gives exactly the minimum if hosts are
restricted to finite graphs. Neither the value $n-1$ nor the asserted
limit is built into the extremum's definition.

\section{Verification and reproduction}
The project pins Lean 4.19.0 and Mathlib v4.19.0, commit
\code{c44e0c8ee63ca166450922a373c7409c5d26b00b}.
The author build and the following strict replays passed:
\begin{quote}\ttfamily\small
lake build\\
lake env lean -DwarningAsError=true Conjecture2187.lean\\
lake env lean -DwarningAsError=true AuthorAudit.lean
\end{quote}
Full reproduction instructions and verification records accompany the package.
The additional audit enumerates every mathematical-module declaration,
including private and generated declarations, and traverses types and
proof values. It rejects owned axioms and unsafe or partial declarations;
only \code{propext}, \code{Classical.choice}, and \code{Quot.sound} may
occur as axioms in the dependency closure. All reasoning is kernel checked,
with no admitted proofs, \code{native\_decide}, custom axioms, or external
numerical computation.

These are local checks and independent review; maintainer acceptance is separate.

\begin{thebibliography}{1}
\small
\bibitem{cfw}
D.~Conlon, J.~Fox, and Y.~Wigderson,
\emph{Three early problems on size Ramsey numbers},
arXiv:2111.05420v2 (2023).
\href{https://arxiv.org/abs/2111.05420v2}{arxiv.org/abs/2111.05420v2}.
Used only for the conventional definition in the abstract.
\end{thebibliography}
\end{document}
